% GENERATED FILE. Edit canonical lessons, then run npm run generate:books.
% canonical-source-hash: fnv1a64:f2073a55b31459d1
% canonical-lessons: UR-C48-khelna, UR-C48-kharidna, UR-C48-bechna, UR-C48-dena, UR-C48-rakhna

\chapter{To play, To buy, To sell, To give, To put}
\label{ch:ur-to-play-to-buy-to-sell-to-give-to-put}
\hlchaptermodality{48}

\begin{chapteropening}
\textbf{By the end of this chapter:} \emph{I can say to play, to buy, to sell, to give and to put, to keep.}

Complete the last lesson of chapter 48: I can say to play, to buy, to sell, to give and to put, to keep.
\end{chapteropening}

\section[\texorpdfstring{\ur{کھیلنا} (khelnā)}{کھیلنا (khelnā)}]{\ur{کھیلنا} (khelnā) — to play}

\label{lesson:UR-C48-khelna}



\begin{quote}
Before the new one: say the Urdu for to walk, then the Urdu for to run.
\end{quote}

\subsection*{You'll want to know: \ur{کھیلنا}}
\textbf{\ur{کھیلنا}} — \emph{khelnā} — \textquotedblleft{}to play\textquotedblright{}.

\begin{cousinweb}
\textbf{\ur{کھیلنا}} (\emph{khelnā}) is inherited from a reconstructed \emph{*khel-}, \textquotedblleft{}play\textquotedblright{} (CDIAL 3918), the verb behind \emph{khel}, \textquotedblleft{}a game\textquotedblright{}. CDIAL links it with Sanskrit \emph{krīḍati}, \textquotedblleft{}plays\textquotedblright{}: if that verb has a non-Aryan source, these forms may share it.
\end{cousinweb}

\begin{grammarlens}[title={What you've built}]
One more word for this chapter.
\end{grammarlens}

\subsection*{Your turn}
\begin{itemize}
\raggedright
  \item \emph{Say it:} \emph{khelnā}
  \item \emph{Say it:} \emph{khelnā}, once more
  \item \emph{Say it:} say \emph{khelnā}
  \item \emph{Recall:} say the Urdu for to walk, then the Urdu for to run, then say \emph{khelnā} again
\end{itemize}

\begin{tcolorbox}[breakable,colback=teal!4,colframe=teal!35!black,title={Before you move on}]
What does \ur{کھیلنا} mean? (\textquotedblleft{}To play\textquotedblright{}.) Say it once more.
\end{tcolorbox}
\section[\texorpdfstring{\ur{خریدنا} (kharīdnā)}{خریدنا (kharīdnā)}]{\ur{خریدنا} (kharīdnā) — to buy}

\label{lesson:UR-C48-kharidna}



\begin{quote}
Before the new one: say the Urdu for to run, then the Urdu for to play.
\end{quote}

\subsection*{You'll want to know: \ur{خریدنا}}
\textbf{\ur{خریدنا}} — \emph{kharīdnā} — \textquotedblleft{}to buy\textquotedblright{}.

\begin{cousinweb}
\textbf{\ur{خریدنا}} (\emph{kharīdnā}) is a loan, usually traced to Persian \emph{kharīdan}, \textquotedblleft{}to buy\textquotedblright{}. Urdu took its stem \emph{kharīd-} and added the verb ending \emph{-nā}.
\end{cousinweb}

\begin{grammarlens}[title={What you've built}]
One more word for this chapter.
\end{grammarlens}

\subsection*{Your turn}
\begin{itemize}
\raggedright
  \item \emph{Say it:} \emph{kharīdnā}
  \item \emph{Say it:} \emph{kharīdnā}, once more
  \item \emph{Say it:} say \emph{kharīdnā}
  \item \emph{Recall:} say the Urdu for to run, then the Urdu for to play, then say \emph{kharīdnā} again
\end{itemize}

\begin{tcolorbox}[breakable,colback=teal!4,colframe=teal!35!black,title={Before you move on}]
What does \ur{خریدنا} mean? (\textquotedblleft{}To buy\textquotedblright{}.) Say it once more.
\end{tcolorbox}
\section[\texorpdfstring{\ur{بیچنا} (bechnā)}{بیچنا (bechnā)}]{\ur{بیچنا} (bechnā) — to sell}

\label{lesson:UR-C48-bechna}



\begin{quote}
Before the new one: say the Urdu for to play, then the Urdu for to buy.
\end{quote}

\subsection*{You'll want to know: \ur{بیچنا}}
\textbf{\ur{بیچنا}} — \emph{bechnā} — \textquotedblleft{}to sell\textquotedblright{}.

\begin{cousinweb}
\textbf{\ur{بیچنا}} (\emph{bechnā}) is a tadbhava, a word inherited and reshaped by sound change. It comes from \emph{*vetyayati}, \textquotedblleft{}makes a price\textquotedblright{}, a form CDIAL reconstructs from its descendants (CDIAL 12100). CDIAL builds it on a word for \textquotedblleft{}hire, price\textquotedblright{}, akin to \emph{vetana}, \textquotedblleft{}wages\textquotedblright{}, and calls other proposals less probable.
\end{cousinweb}

\begin{grammarlens}[title={What you've built}]
One more word for this chapter.
\end{grammarlens}

\subsection*{Your turn}
\begin{itemize}
\raggedright
  \item \emph{Say it:} \emph{bechnā}
  \item \emph{Say it:} \emph{bechnā}, once more
  \item \emph{Say it:} say \emph{bechnā}
  \item \emph{Recall:} say the Urdu for to play, then the Urdu for to buy, then say \emph{bechnā} again
\end{itemize}

\begin{tcolorbox}[breakable,colback=teal!4,colframe=teal!35!black,title={Before you move on}]
What does \ur{بیچنا} mean? (\textquotedblleft{}To sell\textquotedblright{}.) Say it once more.
\end{tcolorbox}
\section[\texorpdfstring{\ur{دینا} (denā)}{دینا (denā)}]{\ur{دینا} (denā) — to give}

\label{lesson:UR-C48-dena}



\begin{quote}
Before the new one: say the Urdu for to buy, then the Urdu for to sell.
\end{quote}

\subsection*{You'll want to know: \ur{دینا}}
\textbf{\ur{دینا}} — \emph{denā} — \textquotedblleft{}to give\textquotedblright{}.

\begin{cousinweb}
\textbf{\ur{دینا}} (\emph{denā}) is a tadbhava, a word inherited from Sanskrit and reshaped by sound change. It comes from \emph{dadāti}, \textquotedblleft{}gives\textquotedblright{} (CDIAL 6141).
\end{cousinweb}

\begin{grammarlens}[title={What you've built}]
One more word for this chapter.
\end{grammarlens}

\subsection*{Your turn}
\begin{itemize}
\raggedright
  \item \emph{Say it:} \emph{denā}
  \item \emph{Say it:} \emph{denā}, once more
  \item \emph{Say it:} say \emph{denā}
  \item \emph{Recall:} say the Urdu for to buy, then the Urdu for to sell, then say \emph{denā} again
\end{itemize}

\begin{tcolorbox}[breakable,colback=teal!4,colframe=teal!35!black,title={Before you move on}]
What does \ur{دینا} mean? (\textquotedblleft{}To give\textquotedblright{}.) Say it once more.
\end{tcolorbox}
\section[\texorpdfstring{\ur{رکھنا} (rakhnā)}{رکھنا (rakhnā)}]{\ur{رکھنا} (rakhnā) — to put, to keep}

\label{lesson:UR-C48-rakhna}



\begin{quote}
Before the new one: say the Urdu for to sell, then the Urdu for to give.
\end{quote}

\subsection*{You'll want to know: \ur{رکھنا}}
\textbf{\ur{رکھنا}} — \emph{rakhnā} — \textquotedblleft{}to put, to keep\textquotedblright{}.

\begin{cousinweb}
\textbf{\ur{رکھنا}} (\emph{rakhnā}) is a tadbhava, a word inherited from Sanskrit and reshaped by sound change. It comes from \emph{rakṣati}, \textquotedblleft{}guards\textquotedblright{} (CDIAL 10547).
\end{cousinweb}

\begin{grammarlens}[title={What you've built}]
One more word for this chapter.
\end{grammarlens}

\subsection*{Your turn}
\begin{itemize}
\raggedright
  \item \emph{Say it:} \emph{rakhnā}
  \item \emph{Say it:} \emph{rakhnā}, once more
  \item \emph{Say it:} say \emph{rakhnā}
  \item \emph{Recall:} say the Urdu for to sell, then the Urdu for to give, then say \emph{rakhnā} again
\end{itemize}

\begin{tcolorbox}[breakable,colback=teal!4,colframe=teal!35!black,title={Before you move on}]
What does \ur{رکھنا} mean? (\textquotedblleft{}To put, to keep\textquotedblright{}.) Say it once more.
\end{tcolorbox}
